\documentclass[11pt]{article}
\usepackage[utf8]{inputenc}
\usepackage[T1]{fontenc}
\usepackage{amsmath,amssymb}
\usepackage{graphicx}
\usepackage[margin=1in]{geometry}
\usepackage{natbib}
\usepackage{hyperref}
\usepackage{booktabs}

\title{Noncanonical Inflammasome Assembly Requires Caspase-11 Catalytic Activity and Intra-molecular Autoprocessing}

\author{Author List}

\date{}

\begin{document}
\maketitle

\begin{abstract}
Inflammatory caspases are cysteine protease zymogens whose activation following infection or cellular damage occurs within supramolecular organizing centers (SMOCs) known as inflammasomes. Inflammasomes are large oligomeric complexes that recruit caspases to undergo proximity-induced autoprocessing, leading to formation of the enzymatically active form that cleaves downstream targets. Binding of bacterial LPS to its cytosolic sensor, caspase-11 (Casp11), promotes Casp11 aggregation within a high molecular weight complex known as the noncanonical inflammasome, where it is activated to cleave gasdermin D and induce pyroptosis. However, the cellular correlates of Casp11 oligomerization and whether Casp11 forms an LPS-induced SMOC within cells remain unknown. Using fluorescently labeled Casp11, we found that LPS transfection of macrophages induced Casp11 speck formation, providing direct evidence that Casp11 forms LPS-induced specks in macrophages. Unexpectedly, we found that catalytic activity was required for Casp11 to form LPS-induced specks in macrophages. Importantly, both catalytic activity and autoprocessing were required for Casp11 speck formation in an ectopic expression system, and Casp11 processing via an exogenous protease was sufficient to induce Casp11 speck formation. These data reveal a previously undescribed role for Casp11 catalytic activity and self-cleavage in noncanonical inflammasome assembly and indicate that Casp11 catalytic activity and autoprocessing occur upstream of, and mediate, Casp11 oligomerization. Our data provide new insight into the molecular requirements for assembly of Casp11 noncanonical inflammasome complexes.
\end{abstract}

\section{Introduction}

The mammalian innate immune system relies on evolutionarily conserved pattern recognition receptors (PRRs) to detect microbe-derived pathogen-associated molecular patterns (PAMPs) and rapidly initiate protective immune responses \citep{janeway1989approaching,medzhitov2002innate}. Such rapid responses are critical for efficient host defense in the face of pathogens that have the capacity to replicate rapidly and suppress or evade immune detection. A common feature of PRR activation is their inducible assembly into higher order oligomeric protein complexes, termed supra-molecular organizing centers (SMOCs), which mediate the effector function of innate immune signaling pathways \citep{kagan2014smocs}.

Inflammatory caspases are zymogens that contain a C-terminal enzymatic domain and an N-terminal caspase activation and recruitment domain (CARD), which promotes SMOC formation through homotypic interactions \citep{martinon2009inflammasome,latz2012inflammasomes}. The C-terminal enzymatic domain comprises large (17--20 kDa) and small (10--12 kDa) subunits separated by an inter-domain linker (IDL), which contains the self-cleavage site that undergoes processing during caspase activation \citep{boatright2003interdimer}.

Caspase-11 (Casp11) serves as the cytosolic sensor for bacterial lipopolysaccharide (LPS), where direct binding of LPS to the Casp11 CARD promotes its oligomerization into a high molecular weight complex known as the noncanonical inflammasome \citep{shi2014cytoplasmic,kayagaki2013noncanonical,hagar2013cytoplasmic}. Activated Casp11 cleaves gasdermin D (GSDMD) to generate the pore-forming N-terminal fragment that executes pyroptosis \citep{kayagaki2015caspase11,aglietti2016gsdmd}. While it is established that Casp11 autoprocessing at the IDL is important for noncanonical inflammasome signaling \citep{lee2018caspase11,ross2018casp11}, the role of catalytic activity in the assembly of the Casp11 inflammasome itself---as distinct from its downstream effector functions---has not been defined.

In this study, we use fluorescently labeled Casp11 constructs to directly visualize noncanonical inflammasome assembly in macrophages and HEK293T cells. We discover that Casp11 catalytic activity and intra-molecular autoprocessing at the IDL are essential for higher-order oligomerization, establishing that these events occur upstream of, rather than as a consequence of, inflammasome assembly.

\section{Results}

\subsection{Caspase-11 catalytic activity is required for cytosolic LPS-induced speck formation}

Upon sensing cytosolic LPS, Casp11 oligomerizes into high molecular weight complexes known as noncanonical inflammasomes \citep{shi2014cytoplasmic}. To better understand the oligomerization dynamics of the Casp11 inflammasome, we generated lentiviral constructs encoding C-terminal fusions of mCherry with full-length Casp11 (Casp11\textsuperscript{WT}-mCherry) and the corresponding C254A mutant lacking catalytic activity, and transduced these constructs into Casp11$^{-/-}$ bone marrow-derived macrophages (BMDMs).

Casp11\textsuperscript{WT}-mCherry, but not Casp11\textsuperscript{C254A}-mCherry, restored dose-dependent GSDMD cleavage and cytotoxicity in Casp11$^{-/-}$ BMDMs, indicating that the tagged mCherry construct retains function (Figure~1B--C). Both constructs exhibited diffuse cytoplasmic staining in the absence of LPS stimulation. Importantly, Casp11\textsuperscript{WT}-mCherry coalesced into large mCherry specks in response to LPS transfection, consistent with the concept that Casp11 assembles into an LPS-induced higher order complex, or SMOC (Figure~1D).

Casp11\textsuperscript{WT}-mCherry exhibited dose-dependent speck formation in response to increasing intracellular LPS levels, paralleling dose-dependent toxicity (Figure~1E). Surprisingly, Casp11\textsuperscript{C254A}-mCherry did not form specks in response to LPS, even at higher doses, despite being expressed at similar or slightly higher levels. These findings imply that CARD-dependent binding of Casp11 to LPS is not sufficient to mediate Casp11 speck formation within cells, and that Casp11 catalytic activity facilitates assembly of the higher order noncanonical inflammasome.

\subsection{Caspase-11 catalytic activity is required for spontaneous oligomerization in HEK293T cells}

To further dissect these mechanisms, we employed an extensively used system for defining the molecular basis for inflammasome assembly in HEK293T cells \citep{kayagaki2015caspase11,shi2014cytoplasmic}. In 293T cells, ectopic Casp11 expression mediates activation of the noncanonical inflammasome and processing of GSDMD.

Casp11\textsuperscript{WT}-mCherry spontaneously assembled into large aggregates reminiscent of ASC-containing specks, whereas Casp11\textsuperscript{C254A}-mCherry remained diffuse within the cytosol (Figure~2A--B), consistent with our findings in macrophages. Moreover, the pan-caspase inhibitor Z-VAD-FMK (zVAD) blocked spontaneous speck assembly and autoprocessing by Casp11\textsuperscript{WT}-mCherry in a dose-dependent manner (Figure~2C--D). Importantly, catalytic activity was not required for Casp11 self-association \textit{per se}, as Casp11\textsuperscript{WT} and Casp11\textsuperscript{C254A} associate both with themselves and each other in reciprocal co-immunoprecipitation studies.

Speck formation also occurred with a Casp11-citrine fusion, but not with a CARD-only-citrine fusion, providing further evidence that the enzymatic domain is important in Casp11 oligomerization. Together, these data imply that catalytic activity contributes to higher order Casp11 oligomerization.

\subsection{Caspase-11 processing at the interdomain linker is necessary but not sufficient for oligomerization}

We next sought to understand the functional role of Casp11 catalytic activity in inflammasome assembly. Since Casp11\textsuperscript{WT}-mCherry spontaneously assembled into specks in 293T cells, which lack known substrates or signaling components for innate immune pathways, we hypothesized that Casp11 autoprocessing might be important in its oligomerization.

Casp11\textsuperscript{D285A}-mCherry, which contained a mutation in the IDL autoprocessing site, displayed impaired GSDMD cleavage and cell death as previously described \citep{lee2018caspase11,ross2018casp11}, and speck formation was significantly reduced compared to Casp11\textsuperscript{WT}-mCherry (Figure~3A--B).

To test whether autoprocessing at the IDL is sufficient to induce oligomerization, we replaced the endogenous Casp11 cleavage sequence with the cleavage sequence for tobacco etch virus (TEV) protease to generate a TEV-cleavable mCherry-tagged Casp11 construct (Casp11\textsuperscript{TEV}-mCherry) (Figure~3C). We also generated a catalytically inactive version (C254A). Co-expression of TEV protease with either WT or C254A Casp11\textsuperscript{TEV}-mCherry in HEK293T cells resulted in dose-dependent Casp11 cleavage (Figure~3D).

TEV protease induced Casp11 speck formation when co-expressed with Casp11\textsuperscript{TEV}-mCherry, demonstrating that inducible TEV-dependent cleavage could rescue oligomerization of the autoprocessing-deficient mutant (Figure~3E--F). Surprisingly, however, TEV protease-induced processing failed to rescue speck formation in the C254A Casp11\textsuperscript{TEV}-mCherry mutant, suggesting that Casp11 catalytic activity functions at an alternate site or that the catalytic residue Cys-254 plays an additional role in inflammasome assembly.

\subsection{Wild-type caspase-11 rescues speck formation of catalytically inactive caspase-11}

Given the ability of TEV-induced cleavage to induce speck formation, we anticipated that co-expression of wild-type unlabeled Casp11 might rescue oligomerization of catalytic mutant Casp11\textsuperscript{C254A}-mCherry. We therefore co-transfected fixed doses of labeled Casp11-mCherry with increasing amounts of unlabeled wild-type Casp11 (Figure~4A).

Expression of unlabeled Casp11 significantly increased speck formation by Casp11\textsuperscript{C254A}-mCherry (Figure~4C--D), demonstrating that wild-type Casp11 could rescue oligomerization of catalytically inactive Casp11-mCherry species. Unexpectedly, co-expression of untagged Casp11 did not lead to significant levels of Casp11\textsuperscript{C254A}-mCherry processing, except at the highest concentrations (Figure~4B), indicating that Casp11 does not induce significant levels of \textit{trans}-processing in the course of oligomerization into specks.

Additional mutation of the IDL autoprocessing site of Casp11\textsuperscript{C254A}-mCherry (generating C254A/D285A double mutant) did not affect its ability to form specks in the presence of co-transfected WT unlabeled Casp11 (Figure~4E--H). These studies demonstrate that both catalytically inactive, as well as inactive-uncleavable, Casp11 can be recruited into oligomeric inflammasome complexes in the presence of catalytically active Casp11, which must undergo intra-molecular processing to induce SMOC formation.

\subsection{Caspase-11 \textit{cis} autoprocessing mediates noncanonical inflammasome assembly}

To directly test whether rescue is driven by \textit{cis} autoprocessing of the catalytically competent Casp11 species, we co-transfected labeled Casp11\textsuperscript{C254A}-mCherry with either catalytically inactive Casp11\textsuperscript{C254A} or non-cleavable Casp11\textsuperscript{D285A} unlabeled Casp11 species (Figure~5A).

Critically, WT unlabeled Casp11 again rescued speck formation by catalytically inactive Casp11-mCherry. However, neither unlabeled Casp11\textsuperscript{C254A} nor Casp11\textsuperscript{D285A} could restore speck formation for Casp11\textsuperscript{C254A}-mCherry (Figure~5C--D). These findings suggest that \textit{cis} autoprocessing of Casp11 results in assembly of a functional oligomeric complex to which the labeled Casp11 reporter mutants can be recruited (Figure~5E).

Altogether, these findings demonstrate a key role for Casp11 catalytic activity and \textit{cis} autoprocessing in oligomerization and higher order assembly of the Casp11 inflammasome. Furthermore, our findings indicate that these activities occur upstream of higher order Casp11 inflammasome assembly, rather than occurring as a terminal consequence of inflammasome assembly.

\section{Discussion}

In this study, we provide the first direct visualization of noncanonical inflammasome specks in macrophages and reveal an unexpected requirement for caspase-11 catalytic activity and intra-molecular autoprocessing in inflammasome assembly. Our findings fundamentally revise the current model of noncanonical inflammasome activation, in which oligomerization was thought to precede and trigger autoprocessing.

The canonical model of inflammasome assembly posits that upstream sensors recruit caspases into close proximity via homotypic CARD--CARD interactions, and this proximity-induced dimerization then triggers catalytic activation and autoprocessing \citep{lu2014unified}. Our data suggest that for the noncanonical inflammasome, the relationship between oligomerization and catalytic activity is more complex. Specifically, we demonstrate that: (1) LPS-induced Casp11 speck formation in macrophages requires catalytic activity; (2) both catalytic activity and IDL autoprocessing are required for spontaneous Casp11 oligomerization in HEK293T cells; (3) exogenous cleavage at the IDL by TEV protease can rescue speck formation of autoprocessing-deficient but catalytically active Casp11; (4) wild-type Casp11 can rescue speck formation of catalytically inactive Casp11 through \textit{cis} rather than \textit{trans} processing; and (5) the catalytic cysteine plays additional roles beyond IDL cleavage in promoting oligomerization.

These observations support a model in which initial low-level Casp11 dimerization upon LPS binding triggers \textit{cis} autoprocessing at the IDL, and this processed form then nucleates higher-order oligomeric SMOC assembly. This is distinct from the ASC-dependent canonical inflammasome, where the adaptor ASC drives filamentous polymerization independent of caspase catalytic activity \citep{lu2014unified}.

Our finding that \textit{trans}-processing does not significantly contribute to Casp11 inflammasome assembly is particularly noteworthy. While \textit{trans}-cleavage by TEV protease could rescue oligomerization, endogenous Casp11 did not efficiently process catalytic mutant Casp11 in \textit{trans} during the course of speck formation. This suggests that the initial autoprocessing event is predominantly intramolecular, consistent with a model in which LPS binding induces a conformational change that enables \textit{cis} cleavage prior to higher-order oligomerization.

\section{Materials and Methods}

\subsection{Cell culture}
HEK293T cells were maintained in complete DMEM supplemented with 10\% FBS, 10~mM HEPES, 10~mM sodium pyruvate, and 1\% penicillin/streptomycin at 37$^\circ$C with 5\% CO$_2$. Murine BMDMs were differentiated from bone marrow progenitors of C57BL/6 or Casp11$^{-/-}$ mice using 30\% L929-conditioned complete DMEM for 7 days \citep{bjanes2021genetic}.

\subsection{Cloning and mutagenesis}
The Casp11 coding sequence was fused to mCherry on its C-terminus. Catalytic (C254A) and cleavage (D285A) mutants were generated by site-directed mutagenesis (Q5 SDM Kit, New England BioLabs). TEV-cleavable constructs were generated by engineering the TEV protease consensus cleavage sequence (ENLYFQ/G) immediately upstream of cleavage-deficient (D285A) Casp11-mCherry.

\subsection{Transfection and imaging}
Mammalian expression plasmids were transfected into HEK293T cells using Lipofectamine 2000. Cells were imaged by fluorescence microscopy 18 hours post-transfection. Speck formation was quantified from 4 image frames (100--150 cells) per well and 3 wells per condition.

\subsection{LPS transfection of macrophages}
Casp11$^{-/-}$ BMDMs transduced with indicated constructs were primed with Pam3CSK4 for 4~h, followed by transfection with LPS. Pyroptosis was assayed by LDH release 16~h later.

\subsection{Western blotting}
Whole cell lysates were immunoblotted for mCherry, GSDMD, V5-tag, and $\beta$-actin as loading control.

\subsection{Statistical analysis}
All error bars represent mean $\pm$ SEM of triplicate wells, representative of 2--3 independent experiments. Data were analyzed by two-way ANOVA with Sidak's multiple comparison test. Dose-response curves were plotted using non-linear regression (Log2; three parameters).

\bibliographystyle{unsrtnat}
\bibliography{references}

\end{document}
